\documentclass[11pt]{article}
\usepackage[a4paper,margin=1in]{geometry}
\usepackage{amsmath,amssymb}
\usepackage{booktabs,longtable,array,calc}
\usepackage{xcolor}
\usepackage[colorlinks=true,allcolors=blue!45!black]{hyperref}
\usepackage{etoolbox}
\AtBeginEnvironment{longtable}{\small}
\AtBeginEnvironment{quote}{\small}
\setlength{\tabcolsep}{4pt}
\providecommand{\tightlist}{\setlength{\itemsep}{0pt}\setlength{\parskip}{0pt}}
\setcounter{secnumdepth}{-1}
\emergencystretch=3em
\newcommand{\tcell}[1]{\parbox[t]{0.95\linewidth}{#1}}
\begin{document}
\title{What the Accounts Settle, and What They Leave Open: Depletion as a Cost of Production, and the Missing Step from a Rate to a Horizon}
\author{Amin Abaee\\[0.35em]
{\small Independent Researcher}\\[0.55em]
{\small\href{https://orcid.org/0000-0002-0019-1842}{ORCID: 0000-0002-0019-1842}}\\[0.3em]
{\small\href{mailto:amin\_abaee@ut.ac.ir}{amin\_abaee@ut.ac.ir}}}
\date{September 17, 2026}
\maketitle

\emph{A commentary accompanying "Typed Flux Ledgers and Depletion Arithmetic: Conservation, Componentwise Diagnostics, and the Semantics of Depletion Horizons" (main text v39; supplementary v10). Nothing here is proved; the arithmetic is quoted from the main text, and the standards statements are cited. Where the main text's labels are used, they are its own.}

\textbf{Track.} Commentary / correspondence. \textbf{Numbering convention.} Assertions in this commentary are lettered B1, B2, \ldots{} and are \emph{claims about the record}, not results; the companion carries no theorem counter, deliberately.

\vspace{0.8em}

\section*{1. The question this commentary answers}\label{1-the-question-this-commentary-answers}

A reviewer of the main text asked, in effect: the accounting standards now record depletion --- does your article do anything the standards do not? The answer is a single sentence, and this commentary is its defence: \textbf{the standards settle how a depletion \emph{flow} is recorded and valued; they never take the step from a flow to a \emph{horizon}, and the article's contribution lives on that step.} Everything below is either an exact statement of what the standards do, an exhibit of what they do not reach, or a note on how to cite the difference without over-claiming.

The main text already carries the technical content: the reclassification exhibit in its Section 1.2, the horizon arithmetic and its bias in Section 10.2 (Remark 33), the aggregate's inability to certify components in Section 10.1 (Proposition 30), and the classification of "time to depletion" quantities in Section 6.5. What is added here is the interface: what a statistical agency that adopts the 2025 revision can and cannot say with the article's apparatus, and what it must still decide for itself.

\vspace{0.8em}

\section*{2. What changed, stated exactly}\label{2-what-changed-stated-exactly}

\textbf{B1.} In the 2025 System of National Accounts, depletion of \emph{non-produced} natural resources is recorded as a cost of production in the current accounts. In the 2008 revision the same quantity was recorded in the account of other changes in the volume of assets and liabilities. Consequently net domestic product is affected by depreciation --- the 2025 terminology for what the 2008 revision called consumption of fixed capital --- \emph{and} by depletion, and so are net national income and net saving.

\textbf{B2.} The 2025 revision takes this from the SEEA Central Framework (2014), which standardised the definition and recording of depletion; the framework treats depletion of natural resources like depreciation of fixed assets, on the ground that both are used up in production.

\textbf{B3.} Two consequences are worth stating because they are routinely mis-remembered. First, the change is classification, not measurement: \textbf{no extraction rate, no physical quantity and no valuation method of the 2008 accounts is replaced.} The same number moves. Second, the effect is not confined to net aggregates in every case: recording depletion as a cost may also touch gross product where natural-resource production is undertaken by government for its own final use and output is measured by the sum-of-costs approach (2025 SNA \S7.137, as cited in the OECD compilation guide; the guide is the source used here, and the paragraph number is reproduced only as a quotation from it --- see Section 6).

\textbf{B4.} Alongside the depletion entry, the 2025 revision adopts the split-asset approach: the value of a resource, and therefore its depletion charge, is attributed between the legal owner and the extractor according to economic ownership shares, with rents to the legal owner including royalties, surtaxes and permits linked to extraction. For biological resources yielding once-only products, the boundary between cultivated (produced) and non-cultivated (non-produced) assets is redrawn, regeneration of biological resources is recorded as gross fixed capital formation, and depletion is treated as a cost of production.

\textbf{B5.} The last item deserves emphasis, because it is where the two systems differ most sharply from the quantity the main text analyses. Under the 2025 revision a fishery recovering after a closure can be recorded as \emph{investment} while its stock grows, and a fishery in decline is charged a \emph{cost}. Both are descriptions of a comparison between extraction and growth. Neither is a statement about when the resource stops supporting the economy. The accounts, on their own terms, do not ask that question.

\vspace{0.8em}

\section*{3. What the standards define is a rate}\label{3-what-the-standards-define-is-a-rate}

\textbf{B6.} The definition of depletion the 2025 revision inherits is \emph{extraction in excess of the resource's growth} --- a comparison of two rates, evaluated per accounting period, in the units of the period. It carries no horizon, no barrier, no scenario, and no probability. The main text's Section 1.2 says this of the same object: "the step from a rate to a horizon is taken nowhere in them, and it is not taken here either."

That sentence is the whole relationship between the two literatures, and it is worth unpacking once, because each of the three steps across it has been taken implicitly by somebody writing about depletion.

\textbf{B7 (the first step, rate to stock).} From a rate one gets a stock-level statement only by assuming the rate: a constant-drift or a declared drift class. The main text's uniform-drift bounds (its Proposition 6, and the finite-exhaustion argument around it) are exactly the conditional form of this step: given a drift bounded in a declared class, exhaustion time lies in a stated interval, and the interval's \emph{width} is part of the result. An accounts table supplies no drift class, so it cannot be read as a bound on anything.

\textbf{B8 (the second step, stock to threshold).} A stock level is not a depletion event until a barrier is named. The main text treats the barrier as a declaration, not a discovery --- its Definition 5 (scenario conditioned hitting time) and its reserve-classification discussion both turn on this. The accounts are indifferent to which barrier a country names; so is the depletion charge.

\textbf{B9 (the third step, threshold to decision).} Whether a horizon is \emph{safe} depends on a required rate of service and a control margin. That is the article's Theorem 24 and its Proposition 41; the standards have no counterpart, and the absence is not an oversight. A national accounting system reports what happened to values and volumes within a period; it is not a solvency test.

The three steps are the article's subject, and none of them is optional or cosmetic: the first needs an assumption, the second needs a boundary, and the third needs a purpose. The standards need none of these because they ask none of those questions.

\vspace{0.8em}

\section*{4. An exhibit of the gap: the aggregate overshoot date}\label{4-an-exhibit-of-the-gap-the-aggregate-overshoot-date}

The one applied construction in the main text that needs no dataset recomputation is worth restating here, because it shows what a rate-based standard cannot see. Write the components of an ecological footprint as demand \(d_i\) against biocapacity \(b_i\), and form

\[ \tau_{\mathrm{agg}}=365\,\frac{\sum_i b_i}{\sum_i d_i}=365\sum_i w_i r_i,\qquad w_i=\frac{d_i}{\sum_j d_j}, \qquad r_i=\frac{b_i}{d_i}. \]

Then \(\tau_{\mathrm{agg}}\) is a demand-weighted mean of the component ratios, so \(\tau_{\mathrm{agg}}\ge\tau_{\min}:=365\min_i r_i\), with equality only where every component of positive weight coincides, and the difference

\[ \Pi_\tau=\tau_{\mathrm{agg}}-\tau_{\min}=365\sum_i w_i\bigl(r_i-r_{\min}\bigr) \]

is a compensation premium: one-signed, exactly decomposable, unbounded in the dispersion of the components (main text, Section 10.2 and Remark 33). On the world totals of the National Footprint and Biocapacity Accounts as tabulated by Lin et al. (2018), the 2022 aggregate ratio is \(0.584\) of biocapacity to demand, so \(\tau_{\mathrm{agg}} = 213\) d; and because the carbon component carries zero biocapacity, the minimum component ratio is \(0\), so \(\Pi_\tau = 213\) d --- the entire aggregate date is compensation. Restricted to the five components of positive biocapacity with renormalised weights, the same arithmetic gives \(\tau_{\mathrm{agg}} = 538\) d against \(\tau_{\min} = 365\) d, a premium of \(173\) d, with the restricted premium running \(547\) d (1961), \(346\) d (1980), \(251\) d (2000) and \(173\) d (2022): it shrinks as the components converge, as the display says it must.

\textbf{B10.} Two features of that exhibit are invisible to a rate-based standard, and they are the reason the article exists. (i) The zero row is a \emph{published convention}: the Guidebook to the National Footprint and Biocapacity Accounts states that no biocapacity figure is computed for carbon uptake, because carbon demand is charged against forest land biocapacity and a separate carbon biocapacity would double count it (Global Footprint Network, 2021, Section 9.1.2). The convention makes \(\tau_{\min}=0\) identically, so the premium equals the whole aggregate date \emph{by construction} rather than by dispersion. A standard that records flows and values has no place to register that a headline quantity is an artefact of a bookkeeping choice. (ii) A depletion charge is additive across components; an overshoot date is not, and the direction in which the aggregate is optimistic is fixed by the weighting. That is a \emph{semantic} property of the aggregate, which is why the main text calls its subject a question of semantics rather than of accuracy.

\textbf{B11.} The same construction explains a public disagreement that is not about arithmetic. An Earth Overshoot Day for 2022 was announced as 28 July --- the 209th day of that year. The 2022 row of a later edition of the accounts gives \(\tau_{\mathrm{agg}} = 213\) d. The four-day difference is vintage and nowcasting, not a dispute about the formula: the accounts are recomputed for every year of the series at each release. Anyone comparing a published overshoot date with a published component table is comparing two vintages, and the main text's practice of pinning every quoted figure to a named release (its Section 6.5, and the vintage record in the supplementary's S5) is the cheapest available protection. It costs nothing but honesty.

\vspace{0.8em}

\section*{5. Two things a standard cannot buy, and a third it should not be asked for}\label{5-two-things-a-standard-cannot-buy-and-a-third-it-should-not-be-asked-}

\textbf{B12 (eligibility).} The 2025 revision leaves the physical boundary of an asset to the classification of the resource, and for the renewable-resource categories it admits, what counts as an asset is drawn by viability under prevailing technology and prices. The main text's Section 1.2 reads that clause correctly: an eligibility statement of exactly the kind its Remark 36 treats as reclassification rather than as a property of the stock. This is not a criticism of the standard. A statistical standard must classify, and classification criteria are chosen, not discovered. It is a boundary on what an adopted standard can certify: a resource's presence on an asset boundary is a statement about the accounts, and its absence is not a statement about the resource.

\textbf{B13 (components).} An aggregate cannot certify its components. The main text proves this in a form any auditor can use --- its Proposition 30 computes, by linear programme, the worst deficit \(\delta_j^{\ast}(z)\) consistent with a published aggregate \(z\) and the published component bounds --- and its Section 10.1 states the consequence: an aggregate can refute or alarm, never clear. A depletion charge recorded in a country's accounts is such an aggregate. It says the resource was used up in producing this year's net product. It does not say which component of the resource base carried the shortfall, and no revision of the standard will make it say so, because the information is not in the aggregate.

\textbf{B14 (a horizon, not asked for).} It would be a category error to ask the accounts for a horizon, and the main text is explicit that the comparison runs one way: the accounts supply an \emph{instance} of the phenomenon it analyses, not a premise for it. Nothing in the article's Sections 2 to 10 depends on any claim made here about the standards, and nothing here depends on the 2025 revision. If the revision is amended, or a country implements it differently, the article is unchanged. That is the sense in which the interface is a one-way mirror, and it is a design property, not a hedge: the article's objects are the ledger's own.

\vspace{0.8em}

\section*{6. Citation hygiene for anyone writing this comparison}\label{6-citation-hygiene-for-anyone-writing-this-comparison}

Three practical rules, each learned from a specific difficulty in this review cycle.

\textbf{B15. Cite no paragraph numbers of the 2025 revision.} Compilation-stage documents disagree about the structure of the text they describe: the OECD guide as published discusses the SEEA in Chapters 2, 35 and 36, while the same statement in the circulation-stage material refers to Chapters 2, 34 and 35. Paragraph references are therefore unstable at exactly the moment commentators most need them, and a wrong paragraph number in a commentary is a correction that will itself need correcting. Cite the standard by edition and by substance; where a compilation guide supplies a paragraph number, reproduce the \emph{guide} as the source of the number, as B3 does.

\textbf{B16. Distinguish endorsement from implementation.} The 2025 revision was endorsed by the United Nations Statistical Commission at its fifty-sixth session in March 2025; the schedule of country implementation, the availability of physical asset accounts, and the completeness of depletion estimates vary by country and by resource. "Adopted" and "in the published accounts of country X" are different claims, and the second generally requires the supplementary-style test of reading the numbers off the source rather than citing its title.

\textbf{B17. Never quote a depletion number as a date.} This is the article's own negative content, and it applies with full force to the new accounts: a depletion charge is a rate, and rate times a year is a value, not a year. If a national statistical office now writes "depletion of forests, 1.4\% of GDP", that sentence is complete. The inference "so we have N years" is not in it.

\vspace{0.8em}

\section*{7. The one open computation, and what it would buy}\label{7-the-one-open-computation-and-what-it-would-buy}

The main text flags a question it does not settle: the treatment of the carbon component. It keeps carbon \emph{demand} in the ratio set, because that is what the published accounts do and the zero biocapacity row is their convention, and it notes that excluding the carbon demand instead is a different object --- the non-carbon components only --- on which \(\tau_{\mathrm{agg}} > 365\) d states that no included component overshoots within the year, not that no overshoot occurs.

\textbf{B18.} What would settle the status of that remark is a computation on the current accounts rather than a discussion: recompute the restricted and unrestricted series through the most recent release, with carbon demand excluded, on the 2025 vintage of the accounts, and report (i) whether the sign of the premium's \emph{trend} survives exclusion, (ii) how the restricted and unrestricted dates separate over the last decade, and (iii) whether the four-day vintage gap of B11 widens or closes. Those three numbers are enough to decide whether this commentary has one exhibit or a fourth section of its own. As it stands it has one, deliberately: a commentary whose exhibit is an unrun computation is a proposal, and proposals belong in the article's own future work.

That is the whole of the difference between half a paper and a paper. The construction is in the main text; the numbers are not.

\vspace{0.8em}

\section*{8. What a reporter should publish}\label{8-what-a-reporter-should-publish}

Four sentences of practice, which is the only place a commentary can be usefully prescriptive.

1. Publish the component table with the rate comparison, not a horizon: extraction and growth, per component, with the period named. The reader who wants a date can then build one and will know what they assumed. 2. If a date is published, publish the direction of its optimism. For any demand-weighted aggregate of component ratios, the aggregate is greater than or equal to the worst component's date, and the gap is computable from the same table (B10). One extra column is all it takes. 3. Name the release and the pull, not the publisher. "Per the 2024 release" is checkable; "per national accounts" is not, once revisions are routine. 4. Say which classification decisions the figure depends on. A depletion charge depends on the asset boundary, on the ownership split, and on which component carries a zero by convention. All three are bookkeeping. Naming them costs nothing and makes the figure durable across revisions.

None of the four requires the article's machinery. All four require the distinction the article draws between a quantity that answers one question and a quantity that is read as answering another.

\vspace{0.8em}

\section*{References}\label{references}

Global Footprint Network, 2021. Working Guidebook to the National Footprint and Biocapacity Accounts, 2021 edition. Global Footprint Network, Oakland, CA.

Lin, D., Hanscom, L., Murthy, A., Galli, A., Evans, M., Neill, E., Mancini, M.S., Martindill, J., Medouar, F.-Z., Huang, S., Wackernagel, M., 2018. Ecological footprint accounting for countries: Updates and results of the National Footprint Accounts, 2012--2018. Resources 7, 58. https://doi.org/10.3390/resources7030058

OECD, 2025. Measuring Natural Resources in the National Accounts: Compilation guidance for the implementation of the 2025 SNA. OECD, Paris.

United Nations, 2014. SEEA Central Framework: 2012 Technical Implementation. Statistical Papers, Series M No. 96. United Nations, New York. (Issued jointly with the European Commission, the International Monetary Fund, the Organisation for Economic Co-operation and Development and the World Bank.)

United Nations, 2025. System of National Accounts 2025. United Nations Statistics Division, New York. Endorsed by the United Nations Statistical Commission at its fifty-sixth session, March 2025. https://unstats.un.org/unsd/nationalaccount/sna2025.asp

\vspace{0.8em}

\section*{Declarations}\label{declarations}

\textbf{Funding.} None declared.

\textbf{Competing interest.} The author declares no competing interest.

\textbf{Relationship to the main text.} This commentary was written as part of the same revision cycle; its standards statements were checked against the sources above and against the OECD compilation guide, and its arithmetic is quoted from the main text, Section 10.2, without recomputation. It cites no paragraph number of the 2025 revision except where the guide itself is the source of the number.

\textbf{Use of AI assistance.} Drafting and iterative review were assisted by GLM (Z.ai), Qwen (Alibaba Cloud) and DeepSeek AI, as in the main text. All statements about the standards are the author's to confirm before submission, for the reason given in B15.
\end{document}
